\documentclass[10pt, conference]{IEEEtran}
\IEEEoverridecommandlockouts

% ─── Packages ──────────────────────────────────────────────────────────────────
\usepackage{times}
\usepackage{amsmath, amssymb}
\usepackage{graphicx}
\usepackage{booktabs}
\usepackage{array}
\usepackage{multirow}
\usepackage{makecell}         % for \makecell multi-line cells
\usepackage{xcolor}
\usepackage{hyperref}
\usepackage{enumitem}
\usepackage{cite}
\usepackage{caption}
\usepackage{float}
\usepackage{adjustbox}
\usepackage{algorithm}
\usepackage{algorithmic}
\usepackage{balance}

\hypersetup{colorlinks=true, linkcolor=black, citecolor=black, urlcolor=blue}
\captionsetup{font=small, labelfont=bf}

% Allow line-breaks inside p{} cells via \makecell
\renewcommand\cellalign{tl}
\renewcommand\theadfont{\bfseries\small}

% ═══════════════════════════════════════════════════════════════════════════════
\begin{document}
% ═══════════════════════════════════════════════════════════════════════════════

% ─── TITLE ─────────────────────────────────────────────────────────────────────
\title{Federated Learning-Based Diabetic Retinopathy Detection\\
Using CNN with SCAFFOLD, Personalization, and Differential Privacy}

% ─── AUTHORS — two-column IEEE format ─────────────────────────────────────────
\author{
  \IEEEauthorblockN{Sarth Narola}
  \IEEEauthorblockA{
    \textit{Department of Computer Science and Engineering}\\
    \textit{Institute of Technology, Nirma University}\\
    Ahmedabad, India\\
    23bce194@nirmauni.ac.in
  }
  \and
  \IEEEauthorblockN{Nevil Anghan}
  \IEEEauthorblockA{
    \textit{Department of Computer Science and Engineering}\\
    \textit{Institute of Technology, Nirma University}\\
    Ahmedabad, India\\
    23bce195@nirmauni.ac.in
  }
}

\maketitle

% ─── ABSTRACT ──────────────────────────────────────────────────────────────────
\begin{abstract}
Diabetic retinopathy (DR) is a leading cause of preventable blindness worldwide.
Early detection is critical, but most deep learning systems require centralizing
patient images, which conflicts with HIPAA and GDPR. This paper proposes a
federated learning (FL) system for five-class DR grading using the Indian Retinal
Image Dataset (IRDRD). The system combines three key components: (1)~SCAFFOLD to
correct client drift on Non-IID hospital data, (2)~local fine-tuning for
per-hospital personalization, and (3)~DP-SGD for formal differential privacy
guarantees. A direct comparison of FedAvg versus SCAFFOLD on the same IRDRD
partitions is conducted. Expected results show faster convergence, higher
accuracy, and full privacy compliance with SCAFFOLD.
\end{abstract}

\begin{IEEEkeywords}
Diabetic Retinopathy, Federated Learning, SCAFFOLD, FedAvg, Differential
Privacy, Personalized FL, CNN, IRDRD, Non-IID
\end{IEEEkeywords}

% ═══════════════════════════════════════════════════════════════════════════════
\section{Introduction}
% ═══════════════════════════════════════════════════════════════════════════════

\subsection{Background of Diabetic Retinopathy}

Diabetes mellitus is a long-term condition affecting over 500 million people
globally. One of its most serious complications is diabetic retinopathy (DR),
which damages the blood vessels in the retina. If untreated, DR leads to
permanent blindness. It is classified into five stages: No DR (Grade~0), Mild
NPDR (Grade~1), Moderate NPDR (Grade~2), Severe NPDR (Grade~3), and
Proliferative DR (Grade~4)~\cite{r10}.

About 103 million people currently have DR, and this number may reach
700 million by 2045~\cite{r3}. India has one of the world's largest diabetic
populations, yet access to eye doctors and retinal imaging equipment in rural
areas is very limited.

\subsection{Importance of Early Detection}

DR has no clear symptoms in early stages, so patients often notice vision loss
only when it is too late for effective treatment. Catching DR early allows
management through laser therapy, anti-VEGF injections, and lifestyle changes.
Manual screening by ophthalmologists is slow and cannot reach the full at-risk
population. CNNs have shown diagnostic accuracy comparable to expert
clinicians~\cite{r1}, making automated screening a practical solution. However,
these models require large centralized datasets, which conflicts with patient
privacy regulations.

\subsection{Role of Federated Learning}

Federated learning (FL) allows multiple hospitals to jointly train a model
without sharing raw patient data~\cite{r5}. Each client trains a local model on
its own data and sends only weight updates to a central server. The server
aggregates these updates into an improved global model, which is sent back to all
clients. This setup satisfies HIPAA and GDPR requirements because no patient
images ever leave the originating hospital~\cite{r8}.

\subsection{Problem Motivation}

The most common FL algorithm, FedAvg, assumes data is distributed identically
across clients (IID). Real hospital data is not IID---different hospitals have
different patient groups, equipment, and practices. This heterogeneity causes
\textit{client drift}, where local models diverge from the global optimum,
slowing convergence~\cite{r5}.

Additionally, most FL-DR papers have not tested on the IRDRD dataset, and only
the work by Hossain et al.~\cite{r12} applies formal differential privacy. No
paper combines SCAFFOLD, personalization, and differential privacy in a single
system. This research fills all three gaps.

% ═══════════════════════════════════════════════════════════════════════════════
\section{Research Objectives}
% ═══════════════════════════════════════════════════════════════════════════════

The specific objectives of this research are as follows:

\begin{enumerate}[leftmargin=*, label=\textbf{O\arabic*.}, itemsep=3pt]

\item \textbf{Design and implement a federated learning framework for five-class
DR grading} using the Indian Retinal Image Dataset (IRDRD) with a custom CNN
architecture, ensuring that no raw patient images are transmitted between client
nodes and the central server at any stage of training.

\item \textbf{Integrate SCAFFOLD into the federated DR detection pipeline} and
empirically compare it against FedAvg on identical Non-IID data partitions
(Dirichlet $\alpha=0.5$) to quantify SCAFFOLD's practical advantage in terms of
convergence speed and final classification accuracy on retinal fundus images.

\item \textbf{Implement and evaluate per-hospital personalization} by freezing
the globally trained convolutional feature extractor and fine-tuning only the
fully connected classification layers at each client node, demonstrating
measurable improvement over the global model on each hospital's local test
distribution.

\item \textbf{Provide formal differential privacy guarantees} by applying DP-SGD
(clipping norm $C=1.0$, noise multiplier $\sigma=1.1$, $\delta=10^{-5}$,
targeting $\epsilon \leq 5$) during both global federated training and local
personalization, and report the privacy--accuracy trade-off relative to
non-private baselines.

\end{enumerate}

% ═══════════════════════════════════════════════════════════════════════════════
\section{Literature Review}
% ═══════════════════════════════════════════════════════════════════════════════

\subsection{Overview of Reviewed Works}

Fifteen papers on FL for DR detection were reviewed, spanning 2020 to 2025.
Karimireddy et al.~\cite{r5} introduced SCAFFOLD to correct client drift on
Non-IID data, while Abidin and Ismail~\cite{r7} achieved 97.87\% accuracy using
FedSGD in one of the first practical FL-DR systems. Nasajpour et al.~\cite{r6}
combined transfer learning with FedAvg and reached 90.07\%.

Recent papers pushed performance further. Mohan et al.~\cite{r10} modified
FedAvg's loss function to reach 98.6\% across five clients, and G.~N.~P and
L.~B~\cite{r14} used FedProx with MobileNet to reach 98.36\% across 20 clients.
Jayalakshmi and Tamilvizhi~\cite{r2} added homomorphic encryption and reached
98.6\% F1 with reduced communication cost.

Formal differential privacy was addressed only by Hossain et al.~\cite{r12},
who applied DP-SGD to FedAvg but reported an accuracy drop to 79.35\%.
Personalization was explored by Bhulakshmi and Rajput~\cite{r3}, reaching
94.66\% via client-level fine-tuning. Vision Transformers entered the FL-DR
domain via Chetoui and Akhloufi~\cite{r4}, achieving AUC up to 0.979. The
theoretical foundation for this work is Karimireddy et al.~\cite{r5}; however,
SCAFFOLD has never been tested on medical imaging, creating a clear research
opportunity. Table~\ref{tab:litreview} summarizes all fifteen papers in
year-ascending order.

% ─── LITERATURE REVIEW TABLE ─────────────────────────────────────────────────
% Spans both IEEE columns (table*) on a normal portrait page. No landscape.
{%
\setlength{\tabcolsep}{3pt}%
\renewcommand{\arraystretch}{1.6}%
\begin{table*}[!t]
\centering
\caption{Summary of Reviewed Literature on Federated Learning for Diabetic Retinopathy Detection}
\label{tab:litreview}
\scriptsize
% Total usable width for table* ~ 17.4cm (IEEE two-column textwidth)
% Cols: 0.75 + 0.38 + 2.6 + 1.6 + 1.45 + 5.3 + 5.3 = 17.38 cm
\begin{tabular}{|>{\centering\bfseries}p{0.75cm}
                |>{\centering}p{0.38cm}
                |p{2.6cm}
                |p{1.6cm}
                |>{\centering}p{1.45cm}
                |p{5.3cm}
                |p{5.3cm}|}
\hline
\textbf{Work} & \textbf{Year} & \textbf{Objective} & \textbf{Technologies} & \textbf{FL Algo.} & \textbf{Pros} & \textbf{Cons} \\
\hline

%── [5] Karimireddy et al. 2020 ──────────────────────────────────────────────
\cite{r5} & 2020
& Fix client drift; converge faster on Non-IID data
& Control variates; PyTorch
& SCAFFOLD
& \parbox[t]{5.3cm}{\raggedright
    Introduced SCAFFOLD algorithm.\newline
    Proven $\mathcal{O}(1/\epsilon^2)$ convergence guarantee.\newline
    Effectively solves client drift on Non-IID data.}
& \parbox[t]{5.3cm}{\raggedright
    Not evaluated on any medical imaging dataset.\newline
    No built-in privacy mechanism.} \\
\hline

%── [6] Nasajpour et al. 2022 ────────────────────────────────────────────────
\cite{r6} & 2022
& Combine transfer learning with FL for DR detection
& AlexNet; TensorFlow
& FedAvg, FedProx
& \parbox[t]{5.3cm}{\raggedright
    First to combine transfer learning with FL for DR.\newline
    Compared FedAvg \& FedProx; reached 90.07\%.}
& \parbox[t]{5.3cm}{\raggedright
    No SCAFFOLD; no Non-IID analysis.\newline
    No differential privacy.} \\
\hline

%── [7] Abidin & Ismail 2022 ─────────────────────────────────────────────────
\cite{r7} & 2022
& FL-based DR detection with privacy preservation
& CNN; PyTorch
& FedSGD
& \parbox[t]{5.3cm}{\raggedright
    One of the first practical FL-DR systems.\newline
    97.87\% accuracy on Kaggle DR dataset.}
& \parbox[t]{5.3cm}{\raggedright
    No Non-IID handling.\newline
    No advanced FL algorithm or differential privacy.} \\
\hline

%── [4] Chetoui & Akhloufi 2023 ──────────────────────────────────────────────
\cite{r4} & 2023
& Use Vision Transformers in FL for DR detection
& ViT; PyTorch
& FedAvg
& \parbox[t]{5.3cm}{\raggedright
    First use of Vision Transformer in FL for DR.\newline
    AUC up to 0.979 across multiple DR datasets.}
& \parbox[t]{5.3cm}{\raggedright
    No personalization; SCAFFOLD not explored.\newline
    Very high compute cost.} \\
\hline

%── [8] Matta et al. 2023 ────────────────────────────────────────────────────
\cite{r8} & 2023
& Test FL in a real multi-center hospital screening network
& EfficientNet-B5; TensorFlow
& FedAvg (weighted / unweighted)
& \parbox[t]{5.3cm}{\raggedright
    Real multi-center FL deployment with 697k images.\newline
    FL nearly matches centralized AUC (0.9317--0.9522).}
& \parbox[t]{5.3cm}{\raggedright
    Required 15-day training time; no advanced FL.\newline
    No differential privacy or personalization.} \\
\hline

%── [10] Mohan et al. 2023 ───────────────────────────────────────────────────
\cite{r10} & 2023
& Federated DR grading with modified cross-entropy loss
& Custom CNN; PyTorch
& FedAvg (modified loss)
& \parbox[t]{5.3cm}{\raggedright
    Achieved 98.6\% accuracy and 99.3\% specificity.\newline
    Modified loss improves class balance across clients.}
& \parbox[t]{5.3cm}{\raggedright
    No personalization; no formal differential privacy.\newline
    Weak Non-IID data handling.} \\
\hline

%── [12] Hossain et al. 2023 ─────────────────────────────────────────────────
\cite{r12} & 2023
& DP-based FL for DR with reduced communication cost
& CNN; Opacus; PyTorch
& FedAvg + DP-SGD
& \parbox[t]{5.3cm}{\raggedright
    First formal DP-SGD applied to FL for DR.\newline
    49\% reduction in communication rounds.}
& \parbox[t]{5.3cm}{\raggedright
    Accuracy drops to 79.35\% when DP is enabled.\newline
    No SCAFFOLD or personalization support.} \\
\hline

%── [13] Singh & Singh 2023 ──────────────────────────────────────────────────
\cite{r13} & 2023
& DenseNet-based FL for DR with improved feature reuse
& DenseNet; PyTorch
& FedAvg
& \parbox[t]{5.3cm}{\raggedright
    Dense skip connections improve feature reuse.\newline
    Custom MPCE loss function enhances accuracy.}
& \parbox[t]{5.3cm}{\raggedright
    No differential privacy; no Non-IID handling.\newline
    No personalization mechanism.} \\
\hline

%── [1] Raj et al. 2024 ──────────────────────────────────────────────────────
\cite{r1} & 2024
& Multi-hospital FL for DR in low-resource settings
& EfficientNetB0; TensorFlow; Web/Mobile
& FedAvg
& \parbox[t]{5.3cm}{\raggedright
    Multi-hospital FL with mobile and web deployment.\newline
    93.21\% accuracy achieved on unseen hospital data.}
& \parbox[t]{5.3cm}{\raggedright
    No personalization or differential privacy.\newline
    No Non-IID analysis; FedAvg only.} \\
\hline

%── [3] Bhulakshmi & Rajput 2024 ─────────────────────────────────────────────
\cite{r3} & 2024
& Build personalized FL model for DR grading
& Multiple CNNs; PyTorch
& FedAvg + Personalization
& \parbox[t]{5.3cm}{\raggedright
    Client-level personalization achieves 94.66\% accuracy.\newline
    Demonstrates benefit of local fine-tuning in FL.}
& \parbox[t]{5.3cm}{\raggedright
    No differential privacy; no SCAFFOLD.\newline
    Communication cost not studied.} \\
\hline

%── [2] Jayalakshmi & Tamilvizhi 2025 ────────────────────────────────────────
\cite{r2} & 2025
& Privacy-preserving FL with encryption and reduced cost
& ShuffleNet+CSPNet; GRU; HE; PyTorch
& FedAvg + IKMC + HE
& \parbox[t]{5.3cm}{\raggedright
    98.6\% F1 with encrypted FL and client clustering.\newline
    Significantly reduces communication cost.}
& \parbox[t]{5.3cm}{\raggedright
    Very complex multi-component pipeline.\newline
    No SCAFFOLD comparison provided.} \\
\hline

%── [9] Mane & Shekapure 2025 ────────────────────────────────────────────────
\cite{r9} & 2025
& 5-class FL-based DR detection with GAN augmentation
& CNN + GAN; PyTorch
& FedAvg (weighted)
& \parbox[t]{5.3cm}{\raggedright
    FL combined with GAN augmentation for 5-class grading.\newline
    Approximately 97\% accuracy reported.}
& \parbox[t]{5.3cm}{\raggedright
    No advanced FL algorithm used.\newline
    No personalization or differential privacy.} \\
\hline

%── [11] Zende & Sonavane 2025 ───────────────────────────────────────────────
\cite{r11} & 2025
& Unified FL + transfer learning + explainability (XAI)
& ResNet+ShuffleNet; Grad-CAM; SHAP
& FedAvg + XAI
& \parbox[t]{5.3cm}{\raggedright
    Integrates XAI explainability into FL for DR.\newline
    Combines transfer learning and FL in one framework.}
& \parbox[t]{5.3cm}{\raggedright
    No experimental results reported in paper.\newline
    No differential privacy or advanced FL algorithm.} \\
\hline

%── [14] G.N.P & L.B 2025 ────────────────────────────────────────────────────
\cite{r14} & 2025
& Lightweight FL-based DR detection for resource-limited settings
& MobileNet; PyTorch
& FedProx + GFA
& \parbox[t]{5.3cm}{\raggedright
    Lightweight MobileNet achieves 98.36\% accuracy.\newline
    Handles Non-IID data effectively using FedProx (20 clients).}
& \parbox[t]{5.3cm}{\raggedright
    No formal differential privacy; no XAI.\newline
    Binary classification only; not 5-class grading.} \\
\hline

%── [15] Bollimuntha et al. 2025 ─────────────────────────────────────────────
\cite{r15} & 2025
& Improve DR classification with DNN over classical ML
& DNN; Scikit-learn
& Centralized only
& \parbox[t]{5.3cm}{\raggedright
    DNN outperforms classical SVM-based ML for DR classification.}
& \parbox[t]{5.3cm}{\raggedright
    No federated learning; no privacy guarantees.\newline
    Not scalable to multi-center hospital settings.} \\
\hline

\end{tabular}
\end{table*}
}

% ═══════════════════════════════════════════════════════════════════════════════
\section{Research Gap}
% ═══════════════════════════════════════════════════════════════════════════════

After reviewing all fifteen papers, six clear gaps were identified.

\textbf{Gap 1 --- No SCAFFOLD for DR Detection:}
SCAFFOLD corrects client drift and outperforms FedAvg on Non-IID
data~\cite{r5}, yet no reviewed paper has applied it to DR detection or any
medical imaging task.

\textbf{Gap 2 --- No IRDRD-Based Federated Study:}
None of the papers use the Indian Retinal Image Dataset (IRDRD). Models trained
on Western datasets may not generalize well to Indian patients, who constitute a
large fraction of the global diabetic population.

\textbf{Gap 3 --- Non-IID Problem Is Largely Ignored:}
Only Karimireddy et al.~\cite{r5} and G.~N.~P and L.~B~\cite{r14} genuinely
address Non-IID data. Most other papers assume IID distributions or apply simple
weighted averaging without analysing data heterogeneity.

\textbf{Gap 4 --- No Combined SCAFFOLD + Personalization + DP:}
Each component appears in isolation: personalization in Bhulakshmi and
Rajput~\cite{r3}, differential privacy in Hossain et al.~\cite{r12}, and
SCAFFOLD in Karimireddy et al.~\cite{r5}. No paper brings all three together
in one system.

\textbf{Gap 5 --- No FedAvg vs.\ SCAFFOLD Comparison on Retinal Data:}
No paper directly compares FedAvg and SCAFFOLD on the same retinal dataset under
Non-IID conditions, making it impossible to quantify SCAFFOLD's practical benefit
for DR screening.

\textbf{Gap 6 --- Five-Class Grading with Privacy Is Rare:}
Papers performing five-class grading~\cite{r9, r10} do not use advanced FL
algorithms or formal privacy mechanisms, limiting their clinical applicability.

% ═══════════════════════════════════════════════════════════════════════════════
\section{Proposed Methodology}
% ═══════════════════════════════════════════════════════════════════════════════

\subsection{Methodology Overview}

The proposed system is an FL framework for five-class DR grading that extends
standard FL with: (1)~SCAFFOLD aggregation to correct client drift on Non-IID
data, (2)~local fine-tuning for per-hospital personalization after global
convergence, and (3)~DP-SGD for formal differential privacy. The local model is
a custom CNN. All experiments run on the IRDRD dataset across five simulated
federated client nodes using Python, PyTorch, and Flower (flwr). Both FedAvg and
SCAFFOLD are compared on identical Non-IID data splits.

\subsection{System Architecture}

The central server holds the global CNN and the SCAFFOLD control variates. Each
client holds a private IRDRD partition and trains locally. Only weight updates
are transmitted---no raw patient images ever leave the client.
Fig.~\ref{fig:architecture} illustrates the architecture.

\begin{figure}[!htbp]
\centering
\fbox{
\begin{minipage}{0.90\columnwidth}
\centering
\vspace{4pt}
\small\textbf{Proposed Federated Architecture}\\[4pt]
\scriptsize
\textbf{Central Server:} Hosts global CNN and SCAFFOLD control variates $c$.
Aggregates $\Delta w_k$ and $\Delta c_k$ from all $N$ clients.
Broadcasts updated model $w^{t+1}$ each round.\\[4pt]
\textbf{Client 1:} Local IRDRD partition $D_1$, model $w_1$, control variate $c_1$.\\
\textbf{Client 2:} Local IRDRD partition $D_2$, model $w_2$, control variate $c_2$.\\
$\vdots$\\
\textbf{Client N:} Local IRDRD partition $D_N$, model $w_N$, control variate $c_N$.\\[4pt]
\textbf{Transmission:} Only $\Delta w_k$ and $\Delta c_k$ are sent.
\textbf{No raw images shared.}
\vspace{4pt}
\end{minipage}
}
\caption{High-level architecture of the proposed federated learning system.}
\label{fig:architecture}
\end{figure}

\subsection{Step-by-Step Workflow}

\begin{enumerate}[leftmargin=*, label=\textbf{Step \arabic*.}, itemsep=2pt]

\item \textbf{Dataset Partitioning.}
IRDRD is split into $N$ non-overlapping client partitions using Dirichlet
sampling ($\alpha=0.5$) to simulate Non-IID hospital data. Each partition is
split 70\% train / 15\% validation / 15\% test.

\item \textbf{Local Preprocessing.}
Each client independently resizes images to $224\times224$, applies CLAHE for
contrast enhancement, normalizes pixel values, and augments training data (flips,
rotation $\pm15^{\circ}$, brightness and contrast jitter).

\item \textbf{Global Model Initialization.}
The server initializes CNN weights with He initialization. For SCAFFOLD, the
global control variate is set to $c^0=0$. Both are broadcast to all $N$ clients.

\item \textbf{Local Training with DP-SGD.}
Each client receives $w^t$ (and $c^t$ for SCAFFOLD), trains for $E=5$ local
epochs using DP-SGD (clipping $\ell_2$ norm $\leq C=1.0$; Gaussian noise
$\mathcal{N}(0,\sigma^2 C^2\mathbf{I})$), and computes local update
$\Delta w_k = w_k^{t,\mathrm{local}}-w^t$. For SCAFFOLD, the corrected gradient
$g_k^{\mathrm{corr}}=g_k-c_k+c^t$ is applied and the local control variate is
updated.

\item \textbf{Update Transmission.}
Each client sends $\Delta w_k$ (and $\Delta c_k$ for SCAFFOLD) to the server.
Gaussian noise ensures DP guarantees under the composition theorem.

\item \textbf{Server Aggregation.}
\textit{FedAvg}: weighted average of client weights proportional to dataset size.
\textit{SCAFFOLD}: weighted average of weight updates applied to global model,
plus control variate aggregation to produce $c^{t+1}$.

\item \textbf{Broadcast.}
Updated global model $w^{t+1}$ (and $c^{t+1}$) is sent to all clients for the
next round.

\item \textbf{Personalization.}
After global convergence each client freezes Blocks 1--4 and GAP, then
fine-tunes only the FC layers for 10 epochs (Adam, $\eta=1\times10^{-5}$,
DP-SGD applied).

\end{enumerate}

\subsection{CNN Model Architecture}

Each client uses a custom CNN for five-class DR classification from
$224\times224$ color fundus images. Table~\ref{tab:cnnarch} summarizes the
architecture.

\begin{table}[!htbp]
\centering
\caption{Proposed CNN Architecture for DR Classification}
\label{tab:cnnarch}
\scriptsize
\begin{adjustbox}{max width=\columnwidth}
\begin{tabular}{llp{1.8cm}p{2.3cm}}
\toprule
\textbf{Layer} & \textbf{Operation} & \textbf{Output Shape} & \textbf{Notes} \\
\midrule
Input   & ---                                      & $224\!\times\!224\!\times\!3$ & RGB retinal image \\
Block 1 & Conv(32)+BN+ReLU+MaxPool                 & $112\!\times\!112\!\times\!32$ & Feature extraction \\
Block 2 & Conv(64)+BN+ReLU+MaxPool                 & $56\!\times\!56\!\times\!64$  & Feature extraction \\
Block 3 & Conv(128)+BN+ReLU+MaxPool                & $28\!\times\!28\!\times\!128$ & Feature extraction \\
Block 4 & Conv(256)+BN+ReLU+MaxPool                & $14\!\times\!14\!\times\!256$ & Frozen in personalization \\
GAP     & Global Average Pooling                   & $256$                        & Frozen in personalization \\
FC1     & Dense(512,ReLU)+Dropout(0.5)             & $512$                        & Fine-tuned \\
FC2     & Dense(256,ReLU)+Dropout(0.3)             & $256$                        & Fine-tuned \\
Output  & Dense(5,Softmax)                         & $5$                          & Fine-tuned \\
\bottomrule
\end{tabular}
\end{adjustbox}
\end{table}

\noindent
Loss: categorical cross-entropy. Optimizer: Adam ($\eta=1\times10^{-4}$). Batch
normalization stabilizes training across heterogeneous client data; Global
Average Pooling reduces overfitting; Dropout (0.5, 0.3) provides regularization
during both federated training and personalization.

\subsection{Federated Learning Algorithms}

\subsubsection{FedAvg (Baseline)}

FedAvg~\cite{r6} is the standard FL baseline. Each round all $N$ clients train
locally for $E$ epochs; the server aggregates weights:
\begin{equation}
  w^{t+1} = \sum_{k=1}^{N} \frac{n_k}{n}\,w_k^{t,\mathrm{local}}
  \label{eq:fedavg}
\end{equation}
where $n_k$ is client $k$'s dataset size. FedAvg is simple but suffers from
client drift under Non-IID data.

\subsubsection{SCAFFOLD (Proposed Primary Algorithm)}

SCAFFOLD, introduced by Karimireddy et al.~\cite{r5}, fixes client drift via
control variates $c_k$ (per client) and $c$ (global). The corrected local SGD
step is:
\begin{equation}
  w_k \leftarrow w_k - \eta_l\!\left(g_k(w_k) - c_k + c\right)
  \label{eq:scaffold_step}
\end{equation}
After local training, client $k$ transmits:
\begin{align}
  \Delta w_k &= w_k^{\mathrm{local}} - w^t \\
  \Delta c_k &= \frac{1}{E\eta_l}\!\left(w^t - w_k^{\mathrm{local}}\right) - c + c_k
\end{align}
The server updates model and control variates:
\begin{align}
  w^{t+1} &= w^t + \eta_s\cdot\frac{1}{N}\sum_{k=1}^{N}\Delta w_k \\
  c^{t+1} &= c^t + \frac{1}{N}\sum_{k=1}^{N}\Delta c_k
\end{align}
SCAFFOLD converges in $\mathcal{O}(1/\epsilon^2)$ rounds on Non-IID data,
versus $\mathcal{O}(1/\epsilon^{2.5})$ for FedAvg.

\subsubsection{Differential Privacy via DP-SGD}

To go beyond basic FL privacy, DP-SGD~\cite{r12} is applied at each client:
\begin{enumerate}[leftmargin=*, itemsep=1pt]
  \item Compute per-sample gradients $g_i(w)$.
  \item \textbf{Clip:} $\tilde{g}_i = g_i\big/\max\!\left(1,\|g_i\|_2/C\right)$
  \item \textbf{Add noise:} $\hat{g} = \tfrac{1}{B}\!\left(\sum_i\tilde{g}_i +
        \mathcal{N}(0,\sigma^2C^2\mathbf{I})\right)$
  \item \textbf{Update weights} using the noisy gradient.
\end{enumerate}
Parameters: $C=1.0$, $\sigma=1.1$, $\delta=10^{-5}$, targeting $\epsilon\leq5$.
DP-SGD is applied during both federated training and personalization.

\subsection{Personalization Strategy}

After global convergence the global CNN is fine-tuned at each client. Blocks
1--4 and GAP are frozen; only the FC layers are updated for 10 epochs (Adam,
$\eta=1\times10^{-5}$, with DP-SGD). This combines global feature learning with
local adaptation without any extra server communication.

\subsection{FedAvg vs.\ SCAFFOLD: Comparison Design}

Both algorithms run on the same IRDRD Non-IID partitions with identical
settings, as summarized in Table~\ref{tab:comparison}.

\begin{table}[!htbp]
\centering
\caption{Experimental Comparison Design: FedAvg vs.\ SCAFFOLD}
\label{tab:comparison}
\scriptsize
\begin{adjustbox}{max width=\columnwidth}
\begin{tabular}{lll}
\toprule
\textbf{Parameter} & \textbf{FedAvg} & \textbf{SCAFFOLD} \\
\midrule
Dataset            & IRDRD (Non-IID)                       & IRDRD (same partitions)  \\
Clients $N$        & 5                                     & 5                        \\
Dirichlet $\alpha$ & 0.5                                   & 0.5                      \\
Local epochs $E$   & 5                                     & 5                        \\
Comm.\ rounds      & 100                                   & 100                      \\
Local LR $\eta_l$  & $1\!\times\!10^{-4}$                  & $1\!\times\!10^{-4}$     \\
Server LR $\eta_s$ & N/A                                   & 1.0                      \\
Batch size         & 32                                    & 32                       \\
DP                 & DP-SGD ($C\!=\!1.0$, $\sigma\!=\!1.1$) & same                   \\
CNN                & Identical                             & Identical                \\
Eval interval      & Every 10 rounds                       & Every 10 rounds          \\
\midrule
\textbf{Metrics}   & \multicolumn{2}{l}{Global accuracy, per-class F1, convergence curve} \\
\bottomrule
\end{tabular}
\end{adjustbox}
\end{table}

\noindent
The hypothesis is that SCAFFOLD will reach 85\% accuracy in fewer rounds and
achieve higher final accuracy after 100 rounds.

\subsection{Dataset: IRDRD}

IRDRD contains high-resolution color fundus photographs annotated by certified
ophthalmologists using the five-class ICDR scale, representing the Indian
clinical population. The dataset is partitioned into $N=5$ non-overlapping
client subsets using Dirichlet sampling ($\alpha=0.5$); each client uses 70\%
train / 15\% validation / 15\% test splits. Weighted random sampling handles
class imbalance. No images leave any client node at any time.

\subsection{Technologies and Implementation Stack}

\begin{table}[!htbp]
\centering
\caption{Technology Stack}
\label{tab:techstack}
\scriptsize
\begin{adjustbox}{max width=\columnwidth}
\begin{tabular}{ll}
\toprule
\textbf{Component} & \textbf{Tool / Version} \\
\midrule
Programming Language  & Python 3.10 \\
Deep Learning         & PyTorch 2.x (CUDA) \\
Federated Learning    & Flower (flwr) -- custom SCAFFOLD strategy \\
Differential Privacy  & Opacus (PyTorch DP library) \\
Image Processing      & OpenCV, Albumentations \\
Data Handling         & NumPy, Pandas \\
Evaluation            & Scikit-learn (accuracy, F1, confusion matrix) \\
Visualization         & Matplotlib, Seaborn \\
Environment           & Jupyter Notebook / VS Code \\
Hardware              & NVIDIA GPU (CUDA); multi-core CPU \\
\bottomrule
\end{tabular}
\end{adjustbox}
\end{table}

\subsection{Privacy Mechanism Summary}

Privacy is enforced at three levels:
\begin{enumerate}[leftmargin=*, itemsep=1pt]
  \item \textbf{FL-level:} Raw images stay at the client; only model updates are
        transmitted.
  \item \textbf{DP-SGD:} Gaussian noise on gradients gives
        $(\epsilon\leq5,\,\delta=10^{-5})$-DP, protecting against inference
        attacks on weight updates.
  \item \textbf{Personalization:} DP-SGD is also applied during fine-tuning so
        no local patient data leaks through personalized layers.
\end{enumerate}

% ═══════════════════════════════════════════════════════════════════════════════
\section{Expected Outcomes}
% ═══════════════════════════════════════════════════════════════════════════════

The following outcomes are expected after completing experiments on IRDRD:

\begin{enumerate}[leftmargin=*, itemsep=2pt]
  \item \textbf{SCAFFOLD converges faster:} SCAFFOLD is expected to reach 85\%
  accuracy in fewer rounds than FedAvg, based on its proven
  $\mathcal{O}(1/\epsilon^2)$ convergence guarantee~\cite{r5}.

  \item \textbf{Higher final accuracy:} After 100 rounds, SCAFFOLD should reach
  88--94\% overall accuracy versus 82--88\% for FedAvg on the same Non-IID
  partitions.

  \item \textbf{Personalization improves local performance:} Each client's
  fine-tuned model should outperform the global model on its own test set by
  2--5\%, especially for clients with skewed class distributions.

  \item \textbf{DP causes a small accuracy trade-off:} DP-SGD is expected to
  reduce accuracy by 3--6 percentage points compared to non-private training,
  consistent with Hossain et al.~\cite{r12} (79.35\% with DP vs.\ 83.05\%
  without).

  \item \textbf{Five-class grading:} Precision, recall, and F1 will be reported
  for all five DR grades. Minority classes (Severe, Proliferative DR) will have
  lower scores due to class imbalance, partly offset by weighted sampling.

  \item \textbf{Privacy-compliant deployment:} No raw images are shared at any
  stage, satisfying HIPAA and GDPR requirements.
\end{enumerate}

\noindent
All figures above are projections based on algorithm theory and literature
trends. Final quantitative results will be reported after completing experiments.

% ═══════════════════════════════════════════════════════════════════════════════
\section{Conclusion}
% ═══════════════════════════════════════════════════════════════════════════════

This paper presented the design of a federated learning system for five-class DR
detection that brings together SCAFFOLD, personalization, and differential
privacy in a single framework. All experiments run on the IRDRD dataset, which
represents Indian patients and has not been used in any prior FL-based DR study.

The literature review confirmed that no existing paper applies SCAFFOLD to
medical imaging, and no paper combines all three proposed components. The gap
analysis demonstrates that this system addresses more research dimensions
simultaneously than any prior work.

The FedAvg vs.\ SCAFFOLD comparison will provide the first empirical evidence of
SCAFFOLD's advantage on retinal data under Non-IID conditions, contributing a
practical and reproducible blueprint for privacy-preserving multi-hospital DR
screening.

% ─── REFERENCES ────────────────────────────────────────────────────────────────
\begin{thebibliography}{15}

\bibitem{r1}
G. M. Raj, M. G. Morley, and M. Eslami, ``Federated Learning for Diabetic
Retinopathy Diagnosis: Enhancing Accuracy and Generalizability in
Under-Resourced Regions,'' \textit{arXiv:2411.00869}, Oct. 2024.

\bibitem{r2}
R. Jayalakshmi and T. Tamilvizhi, ``Privacy Preservation in Diabetic Disease
Prediction Using Federated Learning Based on Efficient Cross Stage Recurrent
Model,'' 2025.

\bibitem{r3}
D. Bhulakshmi and D. S. Rajput, ``FedDL: Personalized Federated Deep Learning
for Enhanced Detection and Classification of Diabetic Retinopathy,''
\textit{PeerJ Computer Science}, vol. 10, e2508, Dec. 2024.

\bibitem{r4}
M. Chetoui and M. A. Akhloufi, ``Federated Learning for Diabetic Retinopathy
Detection Using Vision Transformers,'' 2023.

\bibitem{r5}
S. P. Karimireddy, S. Kale, M. Mohri, S. Reddi, S. Stich, and A. T. Suresh,
``SCAFFOLD: Stochastic Controlled Averaging for Federated Learning,'' in
\textit{Proc. 37th Int. Conf. Machine Learning (ICML)}, 2020.

\bibitem{r6}
M. Nasajpour, M. Karakaya, S. Pouriyeh, and R. M. Parizi, ``Federated Transfer
Learning For Diabetic Retinopathy Detection Using CNN Architectures,'' in
\textit{Proc. IEEE SoutheastCon}, 2022, pp. 655--661.

\bibitem{r7}
N. Z. Abidin and A. R. Ismail, ``Federated Deep Learning for Automated
Detection of Diabetic Retinopathy,'' 2022.

\bibitem{r8}
S. Matta \textit{et al.}, ``Federated Learning for Diabetic Retinopathy
Detection in a Multi-center Fundus Screening Network,'' in \textit{Proc. 45th
Annual Int. Conf. IEEE EMBC}, 2023.

\bibitem{r9}
A. Mane and S. Shekapure, ``Detection and Classification of Diabetic
Retinopathy Using a Federated Learning System Ensuring Data Privacy,'' 2025.

\bibitem{r10}
N. J. Mohan, R. Murugan, T. Goel, and P. Roy, ``DRFL: Federated Learning in
Diabetic Retinopathy Grading Using Fundus Images,'' \textit{IEEE Trans.\ Parallel
Distrib.\ Syst.}, vol. 34, no. 6, pp. 1789--1802, Jun. 2023.

\bibitem{r11}
T. P. Zende and S. P. Sonavane, ``Combining Federated and Transfer Learning
with Explainable AI for Diabetic Retinopathy Detection,'' 2025.

\bibitem{r12}
I. Hossain, S. Puppala, and S. Talukder, ``Collaborative Differentially Private
Federated Learning Framework for the Prediction of Diabetic Retinopathy,''
2023.

\bibitem{r13}
A. Singh and K. K. Singh, ``FedDDR: A Federated Improved DenseNet for
Classification of Diabetic Retinopathy,'' in \textit{Proc. 6th Int. Conf.
Informatics and Data-Driven Medicine (IDDM)}, Nov. 2023.

\bibitem{r14}
G. N. P and L. B, ``DR-FEDPAM: Detection of Diabetic Retinopathy Using
Federated Proximal Averaging Model,'' \textit{J. Electronics, Electromedical
Engineering, and Medical Informatics}, vol. 7, no. 4, pp. 1259--1271,
Oct. 2025.

\bibitem{r15}
M. Bollimuntha \textit{et al.}, ``Enhanced Diabetic Retinopathy Classification
Using Deep Neural Network,'' \textit{Int. J. Research in Engineering and
Science}, vol. 13, no. 5, pp. 1--7, May 2025.

\end{thebibliography}

\end{document}